%% Provenance-Carrying Evidence Fusion for Biomedical Prediction
%%
%% Result tables are generated from run receipts.  Regenerate before rendering:
%%   python3 generate_wm_pln_fusion_tables.py --out generated \
%%       --cafa5-differential <summary.json> \
%%       --biopathnet-candidate <candidate_summary.json> \
%%       --biopathnet-plan <run_plan.json>
%%
%% Render with:
%%   pdflatex wm-pln-provenance-fusion.tex && pdflatex wm-pln-provenance-fusion.tex

\documentclass[11pt]{article}

\usepackage[margin=1in]{geometry}
\usepackage{amsmath}
\usepackage{amssymb}
\usepackage{amsthm}
\usepackage{booktabs}
\usepackage{float}
\usepackage{listings}
\usepackage{url}
\usepackage[hyperindex,breaklinks,hidelinks]{hyperref}
\hypersetup{
  pdftitle={Provenance-Carrying Evidence Fusion for Biomedical Prediction},
  pdfsubject={A typed evidence reducer, and the pre-registration of two experiments}
}

\lstset{basicstyle=\ttfamily\footnotesize,breaklines=true,frame=single}
\restylefloat{table}

\newtheorem{definition}{Definition}
\newtheorem{proposition}[definition]{Proposition}

\newcommand{\code}[1]{\texttt{#1}}
\newcommand{\kind}[1]{\ensuremath{\mathsf{#1}}}

\title{Provenance-Carrying Evidence Fusion for Biomedical Prediction\\
  \large A typed evidence reducer, and the pre-registration of two experiments}
\author{MeTTapedia Project}
\date{15 July 2026}

\begin{document}
\maketitle

\begin{abstract}
Biomedical prediction increasingly proceeds by fusing heterogeneous evidence:
sequence homology, learned embeddings, annotation transfer, and neural link
predictors. The quantities being fused are heterogeneous \emph{in kind} --- some
are ordinal similarities, some are binary observations, few are calibrated
probabilities --- and they are mutually dependent, because they are computed from
overlapping data. Common practice combines them anyway, because they all happen
to be numbers in $[0,1]$.

We describe a platform that refuses to do this. Evidence is typed by kind, and
combinations that are not semantically justified simply have no reduction rule:
noisy-OR is defined only on calibrated probabilities that share a calibration
regime and lie in distinct dependence groups, and evidential revision is defined
only on disjoint provenance. The fusion itself is performed by a streaming
grouped reduction that emits, for each group, a cryptographic certificate of the
evidence it consumed.

This report has two purposes. It documents the platform and the verification it
currently passes, and it \emph{pre-registers} two experiments --- a
partial-knowledge protein function task and an lncRNA--target link prediction
task --- including their hypotheses, splits, and analysis, before their outcomes
are known.

\textbf{We claim no biological performance in this report.} No full evaluation
has been run, and no comparison against a trained control exists. Every result
table below is generated mechanically from a run receipt; where a run has not
been performed, the table says so and reports nothing.
\end{abstract}

\section{Introduction}

A prediction that cannot be interrogated is difficult to act on. In biomedical
settings the consumer of a prediction usually needs to know not only its score
but what supported it, whether the supporting evidence was independent, and how
much of the apparent confidence is an artifact of counting the same observation
twice. Evidence fusion is exactly where these questions are decided, and it is
usually where they are quietly discarded.

Two failure modes recur. The first is \emph{kind confusion}: a BLAST bit-score,
a cosine similarity between protein embeddings, and a binary evidence flag are
all real numbers in $[0,1]$, and are therefore combined by formulas --- noisy-OR
most often --- whose derivation assumes they are calibrated probabilities. The
formula is well-defined on the floats; it is meaningless on the quantities. The
second is \emph{dependence blindness}: predictors that share training
annotations or an embedding are treated as independent witnesses, so agreement
between them inflates confidence when it should not.

Our position is that both failures are type errors, and should be prevented the
way type errors are prevented: by making the illegal combination inexpressible
rather than by warning about it in prose.

\subsection*{Contributions}

\begin{enumerate}
  \item A typed evidence discipline (\S\ref{sec:evidence}) in which ordinal
    scores, binary observations, calibrated probabilities, and evidential
    counts are distinct kinds, and unjustified combinations have no reduction
    rule.
  \item A streaming grouped reduction with per-group evidence certificates
    (\S\ref{sec:reducer}), in ordered and external-merge modes whose exact
    agreement is checked by executable differential tests, with bounded
    source-side input memory under direct cursor ingestion.
  \item The pre-registration of two experiments (\S\ref{sec:cafa5},
    \S\ref{sec:lnc}), with hypotheses and analyses fixed in advance.
  \item Verification evidence for the platform (\S\ref{sec:verification}),
    including exact agreement of the load-bearing reducer with an independent
    reference implementation on a real evidence slice.
\end{enumerate}

\section{Typed evidence}\label{sec:evidence}

We distinguish four kinds of evidence. Let $g$ range over \emph{dependence
groups} --- declared equivalence classes of sources that share data, training
annotations, or a model family --- and let $c$ range over \emph{calibration
regimes}.

\begin{definition}[Evidence kinds]
\begin{align*}
  \kind{Ordinal}(s, v)             &\quad \text{a value $v$ on a named ordinal scale $s$; comparable, not a probability}\\
  \kind{Binary}(\sigma, o)         &\quad \text{an observation $o$ recorded by source $\sigma$}\\
  \kind{Prob}(c, g, \sigma, p)     &\quad \text{$p \in [0,1]$ calibrated under regime $c$, from source $\sigma$ in group $g$}\\
  \kind{PLN}(n^{+}, n^{-}, \tau)   &\quad \text{positive/negative evidence counts with a provenance stamp $\tau$}
\end{align*}
\end{definition}

The discipline is carried by which combinations are \emph{defined}.

\paragraph{Maximum is scale-local.} $\kind{Ordinal}(s,v_1)$ and
$\kind{Ordinal}(s,v_2)$ reduce to $\kind{Ordinal}(s,\max(v_1,v_2))$ only when
the scale $s$ is literally the same symbol. Two ordinal scores on different
scales have no maximum, because their order relations are not commensurable.

\paragraph{Dependence-aware noisy-OR.} For
$\kind{Prob}(c,g_1,\sigma_1,p_1)$ and $\kind{Prob}(c,g_2,\sigma_2,p_2)$ the
reduction
\begin{equation}
  p_1 \oplus p_2 = 1 - (1-p_1)(1-p_2)
  \label{eq:noisyor}
\end{equation}
is defined \emph{only if} the calibration regimes agree and $g_1 \neq g_2$.
Two consequences matter. An ordinal score can never enter
Eq.~\ref{eq:noisyor}, because it is not of kind $\kind{Prob}$ and no rule
mediates the two. And two sources within one dependence group can never be
noisy-OR'd, so declaring sources dependent is a statement with teeth rather
than documentation.

\paragraph{Revision requires disjoint provenance.} $\kind{PLN}$ counts revise by
addition only when their provenance stamps are disjoint, $\tau_1 \cap \tau_2 =
\emptyset$. Overlapping provenance has no revision rule; a policy that resolves
the overlap must be selected explicitly by the caller, never guessed.

The engineering content of this section is that ``no rule'' is the mechanism.
There is no runtime flag to override, and no default that silently applies. A
pipeline that attempts an unjustified fusion does not produce a wrong number; it
produces nothing, and fails.

\section{Streaming grouped reduction and certificates}\label{sec:reducer}

Fusion is a grouped reduction: partition evidence by a key (here, a
protein--term pair and its dependence group), reduce each group, then combine
groups. At biomedical scale the evidence does not fit in memory, so the reducer
must stream its input. Materialized result storage remains a separate
consumer-side cost and is not included in the source-memory claim below.

The reducer exposes one surface, \code{group-fold}, in two modes.

\begin{description}
  \item[\code{ordered}] requires keys in canonical non-decreasing order. It
    holds exactly one active accumulator, releases each group when the key
    advances, and \emph{rejects} key regression rather than silently producing
    fragmented groups. Ordering is over the canonical text of the key, which is
    bytewise rather than numeric --- a distinction that matters for integer
    keys and is benign for identifier keys.
  \item[\code{external-merge}] accepts arbitrary input order. It builds sorted
    runs bounded by a configured size, tags every record with a global
    monotonic ordinal, and performs a stable $k$-way merge that feeds the same
    ordered reducer. Equal keys retain encounter order.
\end{description}

\paragraph{Mode-agreement contract.}
\code{external-merge} totally orders records by (key, ordinal) and feeds the
ordered reducer over the same canonical bytes. The intended contract is that
both modes produce identical groups and identical certificates. This is checked
by executable differential tests rather than proved formally; the formal
obligation is stated in \S\ref{sec:limits}.

\paragraph{Certificates.} Each emitted group carries
$(\text{count}, \text{SHA-256})$ over the canonical byte stream of the items
that entered it, with each item length-prefixed before hashing so that the
digest is collision-resistant across item boundaries. Together with the retained
trail or replayable input, the certificate commits to exactly which serialized
evidence produced the number and how many items were consumed; the digest alone
does not recover or name those items.

\section{Experiment 1: partial-knowledge protein function}\label{sec:cafa5}

\subsection{Task and boundary}

The task is the partial-knowledge protein function prediction setting: known
annotations are frozen at a submission cutoff, and the ground truth consists of
annotations newly acquired in a later growth period. The community evaluator is
the scoring authority; we do not define the metric.

The temporal boundary is a pre-registered constraint, not an aspiration. The
final ground truth may be read \emph{only} by target selection and by the
evaluator. It may not fit calibration, dependence groups, thresholds, source
weights, or any fusion parameter.

\subsection{Status of the released predictors}

Publicly released predictors (sequence homology, embedding similarity,
annotation transfer) are treated as \emph{frozen input features and controls},
not as our result. Reporting their scores would be a reproduction, not an
experiment.

Critically, these released scores are \emph{not} calibrated probabilities: the
homology scores are normalised bit-scores, the embedding scores are cosine
similarities, and the annotation-transfer evidence is binary. Under the
discipline of \S\ref{sec:evidence} they are of kind $\kind{Ordinal}$ and
$\kind{Binary}$, and therefore \emph{cannot} enter Eq.~\ref{eq:noisyor} at all.

\subsection{Pre-registered fusion policy}

The primary policy is the conservative one: all released sources are placed in a
single dependence group and reduced by maximum. This assumes the sources are
mutually dependent, which is the defensible default given that they derive from
overlapping sequence and annotation data.

A calibrated noisy-OR policy is available but gated: it requires a pinned
calibration artifact per source, one shared calibration regime, and an explicit
independence declaration for every cross-group pair. Until a calibration is
actually fitted from eligible pre-cutoff data, this policy is not reportable,
and any result obtained under it is not a result.

\section{Experiment 2: lncRNA--target link prediction}\label{sec:lnc}

\subsection{Systems}

Three systems are trained on identical frozen splits across five seeds:

\begin{enumerate}
  \item a freshly trained neural path-reasoning control (no published
    checkpoint is reused);
  \item a relation-path evidence model, which mines typed relation-path
    templates from the background graph and learns positive/negative evidence
    per template under a Beta-smoothed count-to-truth semantics;
  \item a hybrid, combining the control's logit and the path model's evidence
    log-odds through a non-negative $L_2$-regularised logistic stacker.
\end{enumerate}

\paragraph{A naming caution.} The relation-path evidence model is
\emph{PLN-style} --- it uses evidence counts and a count-to-truth reading --- but
it is presently implemented natively rather than executed in the WM-PLN runtime.
We therefore do not call this lane ``WM-PLN'' and do not claim it as a WM-PLN
result. Grounding it in the runtime is stated as outstanding work in
\S\ref{sec:limits}.

\subsection{Leakage controls}

Three controls are mechanical rather than documentary.

\emph{Self-leakage.} When collecting path features for a query edge, the query
edge and its inverse are removed from the background graph, so the model cannot
observe the answer it is asked for.

\emph{Selection isolation.} During training and model selection the test split
is not merely unread by convention: the selection view substitutes the
validation file in place of the test file and records a receipt attesting the
substitution. Training and selection commands refuse to run unless that receipt
is present and affirms that no real test bytes are staged.

\emph{One-use test.} The test evaluation is guarded by an exclusive-create lock
and a per-job state machine that moves each job from pending to in-progress to
consumed. A second claim fails. The lock is claimed before any test byte is
opened.

\subsection{Pre-registered hypothesis and analysis}

\begin{quote}
\textbf{Primary hypothesis.} The hybrid achieves higher MRR on the held-out test
split than the freshly trained neural control.
\end{quote}

The primary analysis is a paired per-query bootstrap confidence interval for the
MRR difference, with the per-seed distribution reported. Hits@$k$, mean rank,
calibration, the path-model-only lane, provenance coverage, and path ablations
are secondary. Selection uses validation only; the test split is scored exactly
once. \textbf{A non-improvement, or an interval spanning zero, is a valid result
and will be reported as such.} No post-test retuning is permitted.

\subsection{Negative controls}

Negatives are sampled by exact type compatibility, and shortfalls are recorded
rather than filled by broadening the type. A sampled negative is a
\emph{sampled} negative: it is a type-matched pair not present in the
supervision set, not a biologically verified falsehood. Because filtering
sampled negatives against held-out edges would itself constitute leakage, a
sampled negative may coincide with a true held-out edge. Calibration measured
against sampled negatives is therefore calibration against that control
distribution, and is not population calibration.

\section{Platform verification}\label{sec:verification}

The following are engineering results about the platform. They are not
biological results.

\subsection{Reducer equivalence on a real evidence slice}

The load-bearing reducer was run against an independent reference
implementation of the same reduction law over a real evidence slice, under the
conservative maximum policy, with evidence partitioned into shards by a
digest of the protein identifier. Table~\ref{tab:reducer-equivalence} is
generated from the gate's receipt.

% Generated by generate_wm_pln_fusion_tables.py from run receipts.
% Do not edit by hand; edits will be overwritten and are not receipt-backed.
\begin{table}[H]
  \centering
  \small
  \caption{Differential equivalence of the load-bearing reducer against the
  independent reference implementation on a real evidence slice.  Both
  implementations consume the same 92,202 input rows and
  emit byte-identical fused predictions.  Status reported by the gate:
  \texttt{passed}.}
  \label{tab:reducer-equivalence}
  \begin{tabular}{ll}
    \toprule
    Quantity & Value \\
    \midrule
    Primary backend & \texttt{cetta-group-fold} \\
    Primary role & \texttt{load-bearing-reducer} \\
    Reference backend & \texttt{sqlite-differential-oracle} \\
    Fusion policy & \texttt{conservative-max} \\
    \midrule
    Input rows & 92,202 \\
    Evidence packets & 92,202 \\
    Dependence groups & 75,151 \\
    Fused candidates & 75,151 \\
    Exact duplicates collapsed & 0 \\
    Partition shards (active) & 8 (8) \\
    \midrule
    Compared rows & 75,151 \\
    \textbf{Maximum absolute error} & \textbf{0.0} \\
    \midrule
    Prediction digest (primary) & \texttt{2235ca4c4a4c\ldots} \\
    Prediction digest (reference) & \texttt{2235ca4c4a4c\ldots} \\
    Shard manifest digest & \texttt{427d72e55090\ldots} \\
    \bottomrule
  \end{tabular}

  \vspace{0.5em}
  \begin{tabular}{lr}
    \toprule
    Evidence source & Rows \\
    \midrule
    all-evidence & 7,590 \\
    blast-cafa5 & 74,353 \\
    prott5-cosine-5 & 10,259 \\
    \bottomrule
  \end{tabular}
\end{table}


Two facts are worth drawing out. The maximum absolute error between the two
implementations is exactly zero over all compared rows, and the two
implementations emit the same prediction digest --- so the agreement is
byte-level, not merely numerical. The sharded, streaming reducer and a
straightforward relational reference compute the same function on real data.

\subsection{Pre-registration state}

Table~\ref{tab:preregistration} evidences that the frozen plan's test policy
remains unarmed.

% Generated by generate_wm_pln_fusion_tables.py from run receipts.
% Do not edit by hand; edits will be overwritten and are not receipt-backed.
\begin{table}[H]
  \centering
  \small
  \caption{Pre-registration state of the frozen evaluation plan.  The test split
  is reported here solely to evidence that it remains unopened.}
  \label{tab:preregistration}
  \begin{tabular}{ll}
    \toprule
    Quantity & Value \\
    \midrule
    Plan fingerprint & \texttt{10b2a52df68a\ldots} \\
    Frozen jobs & 15 \\
    \textbf{Test policy state} & \textbf{\texttt{unarmed}} \\
    \bottomrule
  \end{tabular}
\end{table}


\section{Preliminary observation, and what is absent}

\subsection{One validation candidate}

Exactly one hyperparameter point of the relation-path evidence model has been
trained and scored on validation. It is reported in
Table~\ref{tab:biopathnet-candidate} because suppressing a weak preliminary
observation would be a selection effect.

% Generated by generate_wm_pln_fusion_tables.py from run receipts.
% Do not edit by hand; edits will be overwritten and are not receipt-backed.
\begin{table}[H]
  \centering
  \small
  \caption{Single unselected validation candidate of the relation-path evidence
  model.  Reported status: \texttt{validation-candidate-complete-not-selected-not-test}.  These are
  validation-only observations for one hyperparameter point; they are not a test
  result, not a selected model, and not a comparison against any control.
  Shortfall policy recorded by the run:
  \emph{use-all-available-and-record; never broaden target type}.}
  \label{tab:biopathnet-candidate}
  \begin{tabular}{ll}
    \toprule
    Quantity & Value \\
    \midrule
    Job & \texttt{pln-path-evidence--seed-1729} \\
    Experiment fingerprint & \texttt{10b2a52df68a\ldots} \\
    Model digest & \texttt{a85556b90309\ldots} \\
    Path templates & 33,656 \\
    Depth / min.\ support / path cap & 3 / 2 / 512 \\
    Evidence prior $(\alpha,\beta)$ & (1, 1) \\
    Seed & 1729 \\
    \midrule
    Supervised positives & 4,867 \\
    Sampled negatives (requested) & 155,744 \\
    Sampled negatives (obtained) & 155,618 \\
    Sampled-negative shortfall & 126 \\
    Positives without exact-type control & 2 \\
    \midrule
    Validation queries & 608 \\
    Candidates per query (mean) & 929.2 \\
    MRR & 0.011627 \\
    Mean rank & 346.97 \\
    Hits@1 / @3 / @10 & 0.0000 / 0.0000 / 0.0016 \\
    Brier vs.\ sampled negatives & 0.043446 \\
    Brier examples & 20,064 \\
    \midrule
    Test split opened & \texttt{false} \\
    Test plan state & \texttt{unarmed} \\
    \bottomrule
  \end{tabular}
\end{table}


The MRR is weak. Against a mean candidate set of roughly $929$ entries, a
random ranking would achieve an MRR of approximately $\ln(929)/929 \approx
0.0074$; the observed $0.0116$ is therefore on the order of $1.6\times$ random.
This is an unselected point in a grid, at one seed, on validation, with no
control to compare against. It supports no claim about the method, and we draw
none. It is recorded so that the eventual grid cannot be read as if this point
had never existed.

\subsection{The results that do not yet exist}

% Generated by generate_wm_pln_fusion_tables.py from run receipts.
% Do not edit by hand; edits will be overwritten and are not receipt-backed.
\begin{center}
\fbox{\parbox{0.9\linewidth}{\textbf{PENDING: full partial-knowledge evaluation.} The frozen full-target CAFA5 evaluation has not been run. No receipt exists for this run, so no value is reported here.}}
\end{center}


% Generated by generate_wm_pln_fusion_tables.py from run receipts.
% Do not edit by hand; edits will be overwritten and are not receipt-backed.
\begin{center}
\fbox{\parbox{0.9\linewidth}{\textbf{PENDING: control / path-evidence / hybrid comparison.} The multi-seed training grid and the one-use test evaluation have not been run, so no comparison against a control exists. No receipt exists for this run, so no value is reported here.}}
\end{center}


\section{What ``provenance'' does and does not mean here}\label{sec:provenance}

Precision here is worth more than enthusiasm.

\textbf{What exists.} Every fused prediction carries a trail naming the source
packets that contributed, their dependence groups, and their scores; each
reduction group carries a recomputable count/digest certificate; the link
prediction lane emits witness paths, contributing input-line digests, learned
template evidence, and a training-manifest digest. Separately, a formal
development of provenance-aware evidential inference exists, covering source and
scope support, disjointness, revision unions, and scoped forgetting.

\textbf{What does not exist.} These artifacts do not yet compose into a single
verified chain
\[
  \text{dataset/split} \to \text{training} \to \text{model/calibration}
  \to \text{inference} \to \text{prediction} \to \text{evaluation receipt},
\]
and we do not yet prove that a runtime certificate denotes the corresponding
formal derivation. The honest term for the present state is therefore
\emph{provenance certificates and audit trails}, not \emph{proven provenance
chains}. We use the former throughout, and the latter nowhere.

\section{Limitations and threats to validity}\label{sec:limits}

\paragraph{Positive-unlabelled ground truth.} Protein function annotation is
sparse and open-world: the absence of an annotation is not evidence of absence.
Precision measured against such a ground truth is a lower bound, and apparent
false positives may be undiscovered true positives.

\paragraph{Calibration is not yet fitted.} The calibrated fusion policy is
plumbed but its calibration has not been learned from eligible data. No
calibrated result is reported here, and none may be reported until the
calibration is fitted, without consulting final labels.

\paragraph{Scale is designed for, not yet demonstrated.} The reducer streams and
spills by construction, and has been exercised at $10^5$ records and at a real
$10^4$-scale evidence slice. The full evaluation is roughly three orders of
magnitude larger. Working memory is bounded per shard and fails loudly rather
than degrading, so a skewed partition is a diagnosable error rather than a
silent one --- but the full-scale claim is an expectation, not a measurement.

\paragraph{Mode agreement is tested, not proved.} The agreement of the ordered
and external-merge modes, the once-counting property of the reduction, and the
admissibility conditions of Eq.~\ref{eq:noisyor} are enforced by executable
tests. Discharging them as formal obligations against the evidential semantics
is outstanding.

\paragraph{The path model is not yet runtime-grounded.} As noted in
\S\ref{sec:lnc}, the relation-path evidence lane is PLN-style but natively
implemented; it is not a WM-PLN execution result.

\section{Reproducibility}

Datasets, the evaluator, and the upstream model are pinned by content digest and
revision; a run fails rather than proceeds if a pin does not verify. Target
selection, shard partitioning, and negative sampling are seeded and
deterministic. Evaluation outputs are hashed semantically --- rows and columns
canonically ordered before digesting --- so that a digest reflects content
rather than file-writing order.

Every number in the result tables of this report is emitted by a generator that
reads run receipts; the generator is included alongside this document. A table
whose receipt is missing renders as an explicit pending notice. This is the same
discipline the paper argues for, applied to the paper.

\begin{thebibliography}{9}

\bibitem{cafa} P.~Radivojac et~al.
  \newblock A large-scale evaluation of computational protein function prediction.
  \newblock \emph{Nature Methods}, 2013.

\bibitem{cafapk} N.~Zhou et~al.
  \newblock The CAFA challenge reports improved protein function prediction and
  new functional annotations for hundreds of genes through experimental
  screens.
  \newblock \emph{Genome Biology}, 2019.

\bibitem{smin} W.~T. Clark and P.~Radivojac.
  \newblock Information-theoretic evaluation of predicted ontological
  annotations.
  \newblock \emph{Bioinformatics}, 2013.

\bibitem{biopathnet} Y.~Yue et~al.
  \newblock Path-based reasoning for biomedical knowledge graphs with
  BioPathNet.
  \newblock \emph{Nature Biomedical Engineering}, 2025.
  \newblock \url{https://doi.org/10.1038/s41551-025-01598-z}

\bibitem{pln} B.~Goertzel, M.~Ikl\'{e}, I.~F. Goertzel, and A.~Heljakka.
  \newblock \emph{Probabilistic Logic Networks}.
  \newblock Springer, 2009.

\bibitem{prott5} A.~Elnaggar et~al.
  \newblock ProtTrans: Toward understanding the language of life through
  self-supervised learning.
  \newblock \emph{IEEE TPAMI}, 2022.

\end{thebibliography}

\end{document}
